% !TeX root = ./../main.tex
%%%%%%%%%%%%%%%%%%%%%%%%%%%%%%%%%%%%%%%%%%%%%%%
%                                             %
% MOJE PAKIETY, USTAWIENIA ITD                %
%                                             %
%%%%%%%%%%%%%%%%%%%%%%%%%%%%%%%%%%%%%%%%%%%%%%%

% Numeracja wierszy na prawym marginesie.
% Przełącznik \LineNumberingtrue/false znajduje się w pliku main.tex.
\ifLineNumbering
  \usepackage[right,mathlines]{lineno}
  \newcommand{\StartLineNumbering}{\linenumbers}
\else
  \newcommand{\StartLineNumbering}{}
\fi

% Kolor konspektów drukowanych przez znaczniki \agentinfo i \agenttask.
\definecolor{agentgray}{gray}{0.4}

% Bibliografia.
% Plik bibliography/biblio.bib (główna baza literatury) jest wczytywany w config/settings.tex.
% Poniżej dołączane są pozostałe bazy.
\bibliography{bibliography/own_publications}

%%%%%%%%%%%%%%%%%%%%%%%%%%%%%%%%%%%%%%%%%%%%%%%%%%%%%%%%%%%%%%%%%%%%%
% listingi i fragmenty kodu źródłowego
\usepackage{listings}
\lstset{%
language=Python,%
commentstyle=\textit,%
identifierstyle=\textsf,%
keywordstyle=\sffamily\bfseries,%
tabsize=3,%
frame=lines,%
numbers=left,%
numberstyle=\tiny,%
numbersep=5pt,%
breaklines=true,%
escapeinside={@*}{*@},%
}

\usepackage{todonotes}
\todostyle{red}{color=red}
\usepackage{longtable}
\usepackage{multirow}
\usepackage{braket}
\usepackage{algorithm}
\usepackage{algpseudocode}
\floatname{algorithm}{Algorytm}
\usetikzlibrary{arrows.meta,calc,positioning,quantikz2}
\usepgfplotslibrary{groupplots}
\pgfplotsset{compat=1.18}
\usepackage{cleveref}

\makeindex[
  title={Skorowidz},
  intoc,
  columns=2
]
\usepackage[toc]{glossaries-extra}
\makenoidxglossaries

%%%%%%%%%%%%%%%%%%%%%%%%%%%%%%%%%%%%%%%%%%%%%%%
%                                             %
% KONIEC MOICH USTAWIEŃ                       %
%                                             %
%%%%%%%%%%%%%%%%%%%%%%%%%%%%%%%%%%%%%%%%%%%%%%%
